% !TeX root = ../教材题库_全部条目_转写工作稿.tex
% Source: XJ-B1-C02-SU-P060-Y1255
% 第2章 / 小结与复习 小结与复习 / 教材 p.60 / PDF 第67页
% OCR未校对：公式与文字均可能有误；未完成数学TeX转写；禁止直接发给学生
% 原题及相邻原解截图（核对用）：../题目原图/XJ-B1-C02-SU-P060-Y1255-01.webp
% 本文件只在转写完成后作为题干使用；OCR 条目尚不是可靠题干。
\begin{bankitem}{第2章 小结与复习 第4题}{第2章 / 小结与复习 小结与复习 / 教材 p.60 / PDF 第67页}{XJ-B1-C02-SU-P060-Y1255}{OCR未校对：公式与文字均可能有误；未完成数学TeX转写；禁止直接发给学生}
\begin{ocrdraft}
[教材p.60 / PDF第67页]
\par
4．从函数观点看一元二次方程和一元二次不等式，可以感悟到数学知识之间
\par
的紧密关联一一一一相等与不等是可以相互转化的，正是这种相互转化，使我们可以
\par
借助函数图象和方程的解来求解不等式．试结合问题，体会解一元二次不等式的
\par
方法与步骤．
\end{ocrdraft}
\end{bankitem}
